\chapter*{卷三\quad 撤离}
\phantomsection\addcontentsline{toc}{chapter}{卷三\quad 撤离}
\markboth{卷三\quad 撤离}{卷三\quad 撤离}

卷三十二片，为全书分量最重之一卷，所记乃公元前二百二十一年前后最后一次清点与封存。莫斯科档案馆第 117 全宗第 2187 号卷宗第 121 叶所录「设施退役清单（共十九项）」，所指即简二十一；俄文档案抄本向来仅摘引第七项，此处给出全文。此卷释文完整度高于他卷，原因不难推知：清点之册乃最后刻成，写毕即封，未曾历经岁焚。

\begin{jian}{简二十}
元年，会计。在籍十一万九千。可去者十一万九千。留者无。

凡不去者，其名不入此册。
\end{jian}
\chuchu{原件残片 57 号，长 22.0 cm，金属。}
\xiaojì{「元年」何指，无从考订。若指某纪年之始，此片当列于他片之前；然卷内次序依内容排定，不以编号定，故置诸卷首。「在籍」与「可去者」两数相等，第三句「留者无」始为结论，第四句则作限定：凡不去者，其名不入此册。此句与简二十八「外工三百二十人，留」直接冲突，构成全书第一处矛盾。另「留者无」三字刻痕较深，与全片他字不类，或为后人补刻。}

\begin{jian}{简二十一}
退役清单，凡十九项。

第一项：居住环六，封。第二项：气路，断。第三项：水路，断。第四项：库，封。第五项：籍，焚。第六项：器，毁。第七项：垂直运输装置（电梯）已拆，支柱于四十米高度截断，截断处标记，勿挖。第八项：界，断。第九项：外舍，平。第十项：垣门，塞。
\end{jian}
\chuchu{原件残片 60 号，长 28.4 cm，为全书最长一片，下端接续处断裂。}
\xiaojì{全文共十九项，此片存前十项，其余九项据残片 61 号补齐，见简二十二。「勿挖」二字为当时俗语，谓所埋之物不可开掘，非警戒后人之辞；二字附于第七项之末，与前后文例不类，当属后加。十九项之中，封闭十项、断绝四项、销毁三项、迁移一项、标记一项，无一项出于安全防范之考量，全数为处置一座不再运转之设施而设。}

\begin{jian}{简二十二}
第十一项：山，封。第十二项：炉，熄。第十三项：药，弃。第十四项：册，收。第十五项：界石，移。第十六项：标，除。第十七项：名，削。第十八项：口，减。第十九项：门，闭。
\end{jian}
\chuchu{原件残片 61 号，长 23.7 cm，金属，字迹较前片为浅。}
\xiaojì{「口，减」谓登记人口数予以削减，并非减少实口；「名，削」与之相接，谓削去名册，二项合观，足见撤离之后城内不留任何一人之名。「药，弃」之药为何物，全书未再出现；简十六「捣者三千」与简十七「人日三合」或指同一物料，然二片俱无宾语，无可确断。「山，封」列于第十一项，为十九项中唯一只封不毁之物：城可毁、路可断，而山独封。}

\begin{jian}{简二十三}
支柱截，深四十米。填以土，夯实。上覆石板三层。石板之上植草，草不活。

草不活者三处，其一在今垣西。守者言：此三处，冬日不冻。
\end{jian}
\chuchu{原件残片 63 号，长 16.4 cm，陶瓷。}
\xiaojì{「草不活」三字独立成句，为他片所未见，至于是地不宜草木抑或当时实况，不可分辨。「草不活者三处，其一在今垣西」，乃全书唯一以「今」纪地者；「今」指书写之时，则此片或非撤离当时所刻，而属事后追记，性质与卷四各片相近，详下编第七章末。末句「此三处，冬日不冻」系引述「守者」之言；「守者」二字在卷三仅此一见，指撤离之后留守之人，与简三十三「吾辈守之三世」相应。\par{}合观全片：四十米深处埋有截断支柱，其上覆土夯实并盖石板三层，然其上草木不生，且三处冬日不冻。以下仅属笔者推测，残简并无明据：三处草不活，或正因其下埋设之构件尚存余热。}

\begin{jian}{简二十四}
装载：人十一万九千。粟，尽。水，尽。气，尽。器，择要者三千件。册，尽。

凡尽者，不留一。
\end{jian}
\chuchu{原件残片 65 号，长 20.1 cm，金属。}
\xiaojì{「人十一万九千」，与简二十相合。「粟尽、水尽、气尽」三语，谓库存全部携走而不留后用，与简二十一「第十二项：炉，熄」相应，可见撤离之后此城不再维持。「册，尽」谓全部名册悉数携走，而简二十二第十七项复言「名，削」，既已全数携走，何以仍须削名，二者显相抵牾。一种可能之调和为所携者乃底册，所削者乃门籍与悬板之名，然此说无直接材料佐证，仅备一说。}

\begin{jian}{简二十五}
凡去者，各给符一枚。符不还。符上刻其名。

无名者不得登。符失者，不得补。
\end{jian}
\chuchu{原件残片 67 号，长 14.2 cm，陶瓷。}
\xiaojì{「符不还」谓符为个人所执之信物，离城后仍归本人；「符上刻其名」则明示一人一符且须具名。「无名者不得登」与前片「名，削」相抵：既已削其名，何以仍得凭名登船，此为全书最大之一处难以弥合之处。若拟调和，唯有一解，即削名在登船之后而给符在登船之前，然残简并无一片支持此时序，故仍以存疑视之。}

\begin{jian}{简二十六}
凡去者，出界则复老，不复入环，亦不复通音书。

已去之籍，唯存其数，不记其后。
\end{jian}
\chuchu{原件残片 69 号，长 18.7 cm，金属，表面有铅笔痕一道，非刻。}
\xiaojì{此片紧承前片给符之制，明定去者离环赴地即重归衰老，既不复返城中，亦不再与月面互通音问；已去之籍仅留登船总口数，后续地面诸事概不入册。与简二十一「第十七项：名，削」及简二十五「符不还」互为表里，足以说明月面行政记录在撤离封门之后即自行断绝，并无任何地面追踪续档回传城内。表面铅笔痕一道系一九六三年编目核对时所留，非古刻。}

\begin{jian}{简二十七}
地热之衰，二百年间，三十一分。始缓，终急。

急则气先绝，水次之，粟又次之。

末年，气路三断三复，吏以告。主者曰：再断，则会计。
\end{jian}
\chuchu{原件残片 71 号，长 15.3 cm，金属。}
\xiaojì{「三十一分」之分谓千分之一，即 3.1\%，与第 117 全宗第 2187 号卷宗所录地热二百年间下降 3.1\% 相合，乃撤离近因之唯一原件依据。「始缓，终急」四字与简十二「同时而计，不同时而尽」合观，足可复原撤离时序：气最先耗竭，故先备气，水次之，粟居末；撤离之所以必在气绝之前，正因气路一断，居者便无从支应至粟尽之日。\par{}末段两句实为全书最紧要之一处因果：末年气路三断三复，司事之吏上禀其事，主者裁断「再断，则会计」，正与简二十「元年，会计」四字遥相呼应，二者所指详下编第七章。「三断三复」揭示该城此前已屡经故障与勉强修葺，直至最后一次主者决定不再修复。至「主者」为何人，三十七片中无一人留名，此人亦复如是。}

\begin{jian}{简二十八}
外工三百二十人，不入装载之数。留。

留者，给三月之粮。
\end{jian}
\chuchu{原件残片 73 号，长 10.8 cm，金属，下端蚀尽。}
\xiaojì{「留」字之后另有五字为他片所无：「留者，给三月之粮」。此五字界定了留下之性质，属于短期留守而非单纯遗弃。若此解成立，则简二十「留者无」并非误记，而仅就「在装载之数内者」而言；外工本不入装载之数，自不在「留者无」之列。此说虽可调和两片，然残简未明言，仍属推论；三月之粮尽后该三百二十人结局如何，全书亦无一字交代。}

\begin{jian}{简二十九}
六批之数：一批六千，二批八千，三批一万，四批一万三千，五批一万八千，六批六万五千。凡十二万。
\end{jian}
\chuchu{原件残片 75 号，长 25.6 cm，金属，字口深，为他片之最。}
\xiaojì{六批之数为全书最完整之一组数据，与第 117 全宗第 2187 号卷宗及《Лунная прародина》所复原者相合。「凡十二万」与简二十、简二十四所记「十一万九千」相差一千，或系取整，或系统计口径各异；若以六数实加，正得十二万整，则「十一万九千」当别有来源，此一千之差额于下编第六章专门处理。}

\begin{jian}{简三十}
六批者，非一时之事。前五批，岁计之内。第六批，岁计之外。

岁计之内者，有次第；岁计之外者，无次第。
\end{jian}
\chuchu{原件残片 77 号，长 17.9 cm，陶瓷。}
\xiaojì{「岁计之内」谓前五批按常规既定计划执行，「岁计之外」谓第六批不在原定计划之内。此片直接印证第 2187 号卷宗之记载：公元前二百二十一年之第六批乃末批，亦为唯一仓促撤离之一批。「有次第」与「无次第」为本片最要紧之区分：前五批各有先后、名册与配额，第六批则仅存「可去者十一万九千」一言而已。}

\begin{jian}{简三十一}
门闭之日，无送者，无哭声。闭者七人，回首一次。

道上无迹，有迹者风。
\end{jian}
\chuchu{原件残片 79 号，长 12.9 cm，金属，为卷三最后一片。}
\xiaojì{「回首一次」四字为全书唯一带动作描写之记述，与他片纯然记数之文例不类，此片是否与撤离诸片同时同批所刻，尚存疑义。本片未尝提及任何人名，「无送者」三字究指何人无从判定，本书亦不据此立论；末句「有迹者风」四字辞气亦异于他片，疑为后世附加。}
